% ---
% id: meat-core/frank-wolfe/rem:fo-nopert
% kind: rem
% title: "A perturbed gradient bounds another function"
% summary: >-
%   A gradient taken at a perturbed operator is the gradient of a different
%   function, and no statement here relates the two, so the explicit
%   perturbation route has no accounted instance. Every path here
%   differentiates on a face at maximal support; an iterate with a perturbed
%   gradient adds no bracket.
% prominence: standard
% status: accepted
% uses: [meat-core/thm:gradient, thm:firstorder]
% mentions: [meat-core/req:split-form-gradient, meat-core/req:gradient-at-maximal-support, req:first-order-withdrawal]
% assumes: []
% refs: []
% provenance: newly-derived
% lean:
% ---

\begin{remark}\label{rem:fo-nopert}
The rules drawn from Theorem~\ref{thm:gradient}\,(i) and~(ii), the face
rule of Requirement~\ref{req:split-form-gradient} and the one-sided
rule of Requirement~\ref{req:gradient-at-maximal-support}, are the rule of
every path this module specifies, without exception, and there is no
second path: no numerical lower-bound certificate is specified here, so no
route perturbs into the strictly positive cone instead of differentiating
on a face, and neither rule is another path's concern. Requirement~\ref{req:gradient-at-maximal-support} names
two routes to a usable gradient: keep the support maximal, ``for instance
after an explicit perturbation''. Only the first is available to
Theorem~\ref{thm:firstorder}, and (H-NOPERT) says so; it is the only route
this module supplies an accounting for.

A perturbation is not unsound in principle; a gradient taken at a
perturbed operator is the gradient of a \emph{different function}, and
this module states no relation between the two. Concretely: replacing $X$
by $X^{(\delta)}=(1-\delta)X+\delta\,\sigma$ for some $\sigma\succ0$ and
evaluating $\nabla T$ there yields the linearisation of an objective
$O^{(\delta)}_{\alpha}$ built from the perturbed branch operators. To
conclude anything about $\hlow=\inf O^{J}_{\alpha}$ from a bound on
$\inf O^{(\delta)}_{\alpha}$ one needs a \emph{domination} statement,
$O^{(\delta)}_{\alpha}\le O^{J}_{\alpha}$ pointwise or an explicit
two-sided error, and this module proves none. The phrase ``after an
explicit perturbation'' in
Requirement~\ref{req:gradient-at-maximal-support} therefore has no
accounted instance anywhere in this module. It stands there as a statement
about what a solver may do to keep a gradient defined, which is true and
is all it says.

The consequence for a run is a \emph{withdrawal}, not a weakening: an
iterate whose gradient was taken after a perturbation contributes no
bracket to $\hlb$ (Requirement~\ref{req:first-order-withdrawal}), so a run
that perturbs the gradient at every iterate reports no $\hlb$ at all. Reporting a perturbed
bracket as a \tref{first-order} bound would be reporting a bound on a
function nobody asked about.
\end{remark}
